% GENERATED FILE. Edit canonical lessons, then run npm run generate:books.
% canonical-source-hash: fnv1a64:8d5599741c9a4844
% canonical-lessons: LA-C363-concedo, LA-C363-denego, LA-C363-confiteor, LA-C363-obicio, LA-C363-contendo

\chapter{Granting and Denying}
\label{ch:la-granting-and-denying}
\hlchaptermodality{363}

\begin{chapteropening}
\textbf{By the end of this chapter:} \emph{I can grant a point, deny or refuse something, own up to a fault, raise a reproach and keep insisting.}

Complete the last lesson of chapter 363: I can grant a point, deny or refuse something, own up to a fault, raise a reproach and keep insisting.
\end{chapteropening}

\section[concēdō, concēdere]{concēdō, concēdere — to grant, to concede; to give way}

\label{lesson:LA-C363-concedo}



\begin{quote}
Before the new one: say the Latin for to prove wrong, then the Latin for a guess.
\end{quote}

\subsection*{You'll want to know: concēdō, concēdere}
\textbf{concēdō, concēdere} — \textquotedblleft{}to grant, to concede; to give way\textquotedblright{}.

\begin{quote}
Hoc tibi concēdō, sed nōn omnia.
\end{quote}

\begin{quote}
\emph{I grant you this, but not everything.}
\end{quote}

\emph{concēdō, concēdere, concessī, concessum}. Dative of the person, accusative of the point: \emph{hoc tibi concēdō}. English \emph{concede} and \emph{concession}.

\begin{cousinweb}
\emph{concēdō} is built from \emph{con-} and \emph{cēdō} \textquotedblleft{}to go, to yield\textquotedblright{} (Lewis \& Short): to step aside for someone. English \emph{concede} comes from it.
\end{cousinweb}

\begin{grammarlens}[title={What you've built}]
One more word for telling, explaining and describing at B1.
\end{grammarlens}

\subsection*{Your turn}
\begin{itemize}
\raggedright
  \item \emph{Say it:} \emph{concēdō, concēdere}
  \item \emph{Say it:} \emph{concēdō, concēdere}, once more
  \item \emph{Say it:} a sentence of your own with \emph{concēdō, concēdere}
  \item \emph{Recall:} say the Latin for to prove wrong, then the Latin for a guess, then say \emph{concēdō, concēdere} again
\end{itemize}

\begin{tcolorbox}[breakable,colback=teal!4,colframe=teal!35!black,title={Before you move on}]
What does concēdō, concēdere mean? (\textquotedblleft{}To grant, to concede; to give way\textquotedblright{}.) Say it once more.
\end{tcolorbox}
\section[dēnegō, dēnegāre]{dēnegō, dēnegāre — to deny flatly; to refuse}

\label{lesson:LA-C363-denego}



\begin{quote}
Before the new one: say the Latin for a guess, then the Latin for to grant.
\end{quote}

\subsection*{You'll want to know: dēnegō, dēnegāre}
\textbf{dēnegō, dēnegāre} — \textquotedblleft{}to deny flatly; to refuse\textquotedblright{}.

\begin{quote}
Pater mihi pecūniam dēnegāvit.
\end{quote}

\begin{quote}
\emph{Father refused me the money.}
\end{quote}

\emph{dēnegō, dēnegāre, dēnegāvī, dēnegātum}. Stronger than \emph{negō}: a firm no. English \emph{deny} comes through French from \emph{dēnegāre}.

\begin{cousinweb}
\emph{dēnegō} is built from \emph{dē-} and \emph{negō} \textquotedblleft{}to say no\textquotedblright{} (EtymDB; Lewis \& Short). English \emph{deny} comes from it through Old French.
\end{cousinweb}

\begin{grammarlens}[title={What you've built}]
One more word for telling, explaining and describing at B1.
\end{grammarlens}

\subsection*{Your turn}
\begin{itemize}
\raggedright
  \item \emph{Say it:} \emph{dēnegō, dēnegāre}
  \item \emph{Say it:} \emph{dēnegō, dēnegāre}, once more
  \item \emph{Say it:} a sentence of your own with \emph{dēnegō, dēnegāre}
  \item \emph{Recall:} say the Latin for a guess, then the Latin for to grant, then say \emph{dēnegō, dēnegāre} again
\end{itemize}

\begin{tcolorbox}[breakable,colback=teal!4,colframe=teal!35!black,title={Before you move on}]
What does dēnegō, dēnegāre mean? (\textquotedblleft{}To deny flatly; to refuse\textquotedblright{}.) Say it once more.
\end{tcolorbox}
\section[cōnfiteor, cōnfitērī]{cōnfiteor, cōnfitērī — to confess, to admit}

\label{lesson:LA-C363-confiteor}



\begin{quote}
Before the new one: say the Latin for to grant, then the Latin for to deny flatly.
\end{quote}

\subsection*{You'll want to know: cōnfiteor, cōnfitērī}
\textbf{cōnfiteor, cōnfitērī} — \textquotedblleft{}to confess, to admit\textquotedblright{}.

\begin{quote}
Culpam meam cōnfiteor.
\end{quote}

\begin{quote}
\emph{I admit my fault.}
\end{quote}

Deponent: \emph{cōnfiteor, cōnfitērī, cōnfessus sum}. A fuller form of \emph{fateor}, to admit. English \emph{confess}.

\begin{cousinweb}
\emph{cōnfiteor} is built from \emph{con-} and \emph{fateor} \textquotedblleft{}to admit\textquotedblright{} (EtymDB; Lewis \& Short), the vowel weakened inside the compound. English \emph{confess} comes from its participle \emph{cōnfessus}.
\end{cousinweb}

\begin{grammarlens}[title={What you've built}]
One more word for telling, explaining and describing at B1.
\end{grammarlens}

\subsection*{Your turn}
\begin{itemize}
\raggedright
  \item \emph{Say it:} \emph{cōnfiteor, cōnfitērī}
  \item \emph{Say it:} \emph{cōnfiteor, cōnfitērī}, once more
  \item \emph{Say it:} a sentence of your own with \emph{cōnfiteor, cōnfitērī}
  \item \emph{Recall:} say the Latin for to grant, then the Latin for to deny flatly, then say \emph{cōnfiteor, cōnfitērī} again
\end{itemize}

\begin{tcolorbox}[breakable,colback=teal!4,colframe=teal!35!black,title={Before you move on}]
What does cōnfiteor, cōnfitērī mean? (\textquotedblleft{}To confess, to admit\textquotedblright{}.) Say it once more.
\end{tcolorbox}
\section[obiciō, obicere]{obiciō, obicere — to throw in the way; to throw in someone's face}

\label{lesson:LA-C363-obicio}



\begin{quote}
Before the new one: say the Latin for to deny flatly, then the Latin for to confess.
\end{quote}

\subsection*{You'll want to know: obiciō, obicere}
\textbf{obiciō, obicere} — \textquotedblleft{}to throw in the way; to throw in someone's face\textquotedblright{}.

\begin{quote}
Multī eī stultitiam obiciunt.
\end{quote}

\begin{quote}
\emph{Many people reproach him with folly.}
\end{quote}

\emph{obiciō, obicere, obiēcī, obiectum}: dative of the person, accusative of the charge. English \emph{object} and \emph{objection}.

\begin{cousinweb}
\emph{obiciō} is built from \emph{ob} \textquotedblleft{}in the way\textquotedblright{} and \emph{iaciō} \textquotedblleft{}to throw\textquotedblright{} (EtymDB; Lewis \& Short). English \emph{object} comes from its participle \emph{obiectum}.
\end{cousinweb}

\begin{grammarlens}[title={What you've built}]
One more word for telling, explaining and describing at B1.
\end{grammarlens}

\subsection*{Your turn}
\begin{itemize}
\raggedright
  \item \emph{Say it:} \emph{obiciō, obicere}
  \item \emph{Say it:} \emph{obiciō, obicere}, once more
  \item \emph{Say it:} a sentence of your own with \emph{obiciō, obicere}
  \item \emph{Recall:} say the Latin for to deny flatly, then the Latin for to confess, then say \emph{obiciō, obicere} again
\end{itemize}

\begin{tcolorbox}[breakable,colback=teal!4,colframe=teal!35!black,title={Before you move on}]
What does obiciō, obicere mean? (\textquotedblleft{}To throw in the way; to throw in someone's face\textquotedblright{}.) Say it once more.
\end{tcolorbox}
\section[contendō, contendere]{contendō, contendere — to insist, to maintain; to strive; to hurry}

\label{lesson:LA-C363-contendo}



\begin{quote}
Before the new one: say the Latin for to confess, then the Latin for to throw in the way.
\end{quote}

\subsection*{You'll want to know: contendō, contendere}
\textbf{contendō, contendere} — \textquotedblleft{}to insist, to maintain; to strive; to hurry\textquotedblright{}.

\begin{quote}
Contendit sē innocentem esse.
\end{quote}

\begin{quote}
\emph{He insists that he is innocent.}
\end{quote}

\emph{contendō, contendere, contendī, contentum}. With the accusative and infinitive, to insist; with \emph{ad}, to hurry to a place. English \emph{contend}.

\begin{cousinweb}
\emph{contendō} is built from \emph{con-} and \emph{tendō} \textquotedblleft{}to stretch\textquotedblright{} (EtymDB; Lewis \& Short): to stretch every nerve, for a goal or for a claim. English \emph{contend} comes from it.
\end{cousinweb}

\begin{grammarlens}[title={What you've built}]
One more word for telling, explaining and describing at B1.
\end{grammarlens}

\subsection*{Your turn}
\begin{itemize}
\raggedright
  \item \emph{Say it:} \emph{contendō, contendere}
  \item \emph{Say it:} \emph{contendō, contendere}, once more
  \item \emph{Say it:} a sentence of your own with \emph{contendō, contendere}
  \item \emph{Recall:} say the Latin for to confess, then the Latin for to throw in the way, then say \emph{contendō, contendere} again
\end{itemize}

\begin{tcolorbox}[breakable,colback=teal!4,colframe=teal!35!black,title={Before you move on}]
What does contendō, contendere mean? (\textquotedblleft{}To insist, to maintain; to strive; to hurry\textquotedblright{}.) Say it once more.
\end{tcolorbox}
